\chapter{Converging-Evidence Evaluation, Robustness, and Sensitivity}
\label{chap:robustness-sensitivity-validation-design}
\label{chap:experimental-design}

This chapter is the methodological center for evaluating the two observational testbeds without a causal answer key.  It defines the experiment families, population hierarchy, converging-evidence axes, diagnostic framework, result-admission policy, and reproducibility boundary underlying Chapter~\ref{chap:results}.  The axes are deliberately non-equivalent: each investigates a different failure mode, and their findings cannot be added together as proof of clinical correctness, causal identification, or estimator accuracy.  They supplement rather than replace the project DAG, exposure definition, adjustment assumptions, empirical support, and clinical review.

\section{Converging-Evidence Axes}

\subsection{Construct plausibility}

Construct plausibility asks whether a proposed representation has a recognizable clinical rationale and an inspectable operational definition.  The evidence consists of literature-grounded representation design, the preserved rationales and source links, rule inspection, and the bounded informal ICU-clinician consultation reported by the finished study.  There was no blinded formal clinician-rating study, chart adjudication, inter-rater reliability analysis, or clinical construct-validation study.  This axis can identify an implausible or poorly specified construct, but it cannot validate a proxy as a diagnosis, establish that a graph is correct, or show that the encoded state is causally adequate.

\subsection{Proxy recoverability and mortality-relevant predictive information}

Proxy recoverability, detailed in Chapter~\ref{chap:predictive-modeling-proxy-states}, asks whether the deterministic rule-derived labels correspond to learnable patterns in the irregular source-derived representation.  Performance against those labels measures reproduction of a representation-defined target, not clinical correctness.  A separate mortality-prediction task asks whether proxy representations contain information associated with source-observed in-hospital mortality.  Such prognostic information can characterize a resource, but it remains associational and does not establish causal relevance, diagnostic accuracy, intervention utility, correctness of the proxy, or correctness of the DAG.

\subsection{Cross-estimator stability}

Cross-estimator stability asks whether fitted effect patterns are unique to one estimator under a fixed representation and population.  Three distinct dimensions are retained: sign agreement compares direction, rank agreement compares the ordering of exposures within a resource, and magnitude agreement compares numerical scale.  These dimensions are not interchangeable, and estimators may also target or approximate different quantities.  Agreement supplies evidence that a pattern is not confined to one fitted estimator; it does not establish estimator accuracy, estimand equivalence, identification, or truth.

\subsection{Clinical-literature comparison}

Clinical-literature comparison provides contextual corroboration by relating a CliniCause proxy and its model-estimated mortality contrast to the nearest usable published clinical construct.  Where available, the comparison records the represented clinical construct, operational definition, study population, comparison group, mortality horizon, effect scale, adjustment method, reported mortality contrast, and major comparability limitations.  Candidate comparisons are interpreted only at the level supported by the finished paper and supplement.

External studies can differ in clinical definition, severity, eligibility, data source, calendar period, treatment context, observation window, mortality horizon, matching or adjustment, estimand, empirical support, and residual confounding.  The comparison is therefore a way to identify broad compatibility and discrepancies.  It is not a meta-analysis, calibration against causal truth, external validation of a proxy, or proof that a CliniCause estimate is correct.  Numerical synthesis is reserved for Chapter~\ref{chap:results} and is not introduced in this methods chapter.

\section{Populations and Experiment Families}

\subsection{Predictive experiments}
\label{sec:experimental-design:prediction-experiment-matrix}

The predictive comparison crosses PhysioNet 2012 and MIMIC-III with four learned models: STraTS, GRU, GRU-D, and TCN.  Each model predicts dataset-specific deterministic rule-derived proxy labels, not mortality or the later aggregate.  The workflow includes split-aware loading, multi-label training, validation-based checkpoint selection, test evaluation, and probability/binary prediction export; STraTS additionally supports self-supervised pretraining.  These experiments operationalize the proxy-recoverability axis rather than clinical validation.  Training summaries and prediction exports exist for all eight dataset--model combinations.  Processed-pickle hashes, split identifiers, complete producing commands, checkpoint-to-export mappings, and archive-copy histories remain incomplete, so implementation availability, artifact availability, and complete provenance are not interchangeable claims.

\subsection{Causal experiments and population hierarchy}
\label{sec:experimental-design:causal-experiment-matrix}

The causal design crosses two datasets, three estimators, and two sampling conditions, yielding twelve archived run families.  CausalForestDML is primary, LinearDML secondary, and CausalPFN exploratory.  Original cohorts define the primary analyses.  The robustness condition retains all mortality-positive rows and samples mortality-negative rows before exposure-specific matching and estimation; it therefore changes the empirical population and outcome prevalence and is never pooled with the original cohorts.

Every final run uses a deterministic five-source aggregate: one rule-derived table plus thresholded predictions from STraTS, GRU, GRU-D, and TCN.  The runtime order is graph construction, mixed-source aggregation, mortality prediction, matching, CATE estimation, saved-CATE analysis, and permutations.  CausalPFN completes CATE estimation but its EconML-specific saved-analysis and permutation stages are intentionally skipped, not failed.  Within each run, all prespecified dataset-specific non-chronic proxy exposures are evaluated subject to column, graph-node, binary-value, adjustment, support, and artifact checks; the exposure list is a design commitment rather than a set of causal conclusions.

\begin{table}[htbp]
\centering
\caption{Experiment families, populations, and evidential roles.}
\label{tab:experiment-family-summary}
\small
\begin{tabular}{@{}L{0.20\textwidth}L{0.24\textwidth}L{0.25\textwidth}L{0.22\textwidth}@{}}
\toprule
Family & Design & Primary output & Evidential role \\
\midrule
Prediction & Two datasets $\times$ four learned models & Test metrics and binary/probability exports & Comparison of rule-label learnability; not clinical validation \\
Causal, original & Two datasets $\times$ three estimators & Matching and model-estimated CATE summaries & Primary forest, secondary linear, exploratory PFN analyses \\
Causal, downsampled & Same datasets and estimators after outcome-class sampling & Separate matching and CATE summaries & Robustness population only; no pooling or automatic common estimand \\
\bottomrule
\end{tabular}
\end{table}

\section{Executed Settings and Provenance Status}
\label{sec:experimental-design:executed-settings-provenance}

Table~\ref{tab:executed-settings-provenance} separates what the checked archive supports from what only current source defines and what remains unresolved.  \emph{Artifact-supported} denotes a checked summary, export, run summary, or diagnostic record; \emph{log-supported} denotes an argument recorded in archived logs without a complete immutable producing manifest; \emph{source-defined} denotes implemented behavior that must not be mistaken for a recovered historical setting.  \emph{Unresolved} and \emph{not applicable} retain their literal meanings.  This classification is necessary because a current default, a historical path, and an archived output answer different reproducibility questions.

\begingroup
\small
\singlespacing
\sloppy
\setlength{\tabcolsep}{3pt}
\begin{longtable}{@{}L{0.19\textwidth}L{0.42\textwidth}L{0.30\textwidth}@{}}
\caption{Compact specification of executed or execution-supported settings and unresolved lineage.}
\label{tab:executed-settings-provenance}\\
\toprule
Analysis component & Confirmed or execution-supported settings & Unresolved provenance or interpretation boundary \\
\midrule
\endfirsthead
\toprule
Analysis component & Confirmed or execution-supported settings & Unresolved provenance or interpretation boundary \\
\midrule
\endhead
\bottomrule
\endfoot
Prediction scope and artifacts & \emph{Artifact-supported}: PhysioNet 2012 and MIMIC-III crossed with STraTS, GRU, GRU-D, and TCN; all eight dataset--model combinations have checked training summaries and prediction exports.  The task is multi-label prediction of dataset-specific rule-derived proxy labels. & \emph{Unresolved}: processed-pickle hashes, split identifiers, producing source revision, and archive-copy direction.  Artifact availability does not establish an independently auditable out-of-sample membership. \\
Archived predictive run settings & \emph{Artifact-supported}: checked training summaries identify run \texttt{1o10}, training fraction 0.5, validation and test evaluation, and checkpoint selection by AUPRC plus AUROC for each of the eight learned combinations.  \emph{Log-supported}: archived argument records give seed 2023 and \texttt{predict\_split=all}. & \emph{Unresolved}: the seed is not a split manifest; the raw shell command, immutable producing configuration, and checkpoint-to-export mapping are absent.  \texttt{all} is the recorded export argument, not proof of a particular split membership. \\
Prediction export contract & \emph{Artifact-supported}: each checked export has unique \texttt{ts\_id} rows, one \texttt{<label>\_prob} field and one binary \texttt{<label>} field per active target, valid probability ranges, and binary values.  \emph{Source-defined}: the binary export threshold is 0.5. & \emph{Unresolved}: archived exports are hashable, as are candidate best checkpoints, but no immutable record proves which checkpoint produced each export. \\
STraTS pretraining & \emph{Source-defined}: STraTS alone has a separate self-supervised pretraining path before supervised fine-tuning; GRU, GRU-D, and TCN do not use that path.  Archived STraTS training and checkpoint artifacts exist. & \emph{Unresolved}: pretraining-checkpoint, supervised-checkpoint, and final-export lineage is incomplete; source support does not prove the exact pretraining state used historically. \\
Causal run matrix and hierarchy & \emph{Artifact-supported}: two datasets, CausalForestDML, LinearDML, and CausalPFN, and original/outcome-downsampled conditions yield twelve archived run families with stage records.  Original cohorts are primary; outcome-downsampled populations are separate robustness analyses.  CausalForestDML is primary, LinearDML secondary, and CausalPFN exploratory. & \emph{Unresolved}: numbered final causal configuration files and processed-input hashes are missing.  Recorded historical paths are not portable configurations, and current unnumbered configs are only candidates. \\
Exposure, aggregation, and support & \emph{Artifact-supported}: every archived final run records the same five-source mixed aggregate---one rule-derived table plus binary predictions from STraTS, GRU, GRU-D, and TCN.  \emph{Source-defined}: analysis is exposure specific; graph/column checks govern exposure admission, while matching supplies descriptive support and pair-availability diagnostics before CATE interpretation. & \emph{Unresolved}: voter bytes and hashes are not co-archived per run, and matching does not establish positivity, balance, exchangeability, or a common estimand across sampling conditions. \\
Sensitivity evidence & \emph{Artifact-supported}: CausalForestDML and LinearDML records distinguish fitting-time availability, saved-training direct values, later artifact analysis, labelled fallback reconstruction, benchmark status, partial rows, and failures.  \emph{Not applicable}: the archived CausalPFN branch has no equivalent EconML saved-analysis family. & \emph{Unresolved}: evidence classes are not numerically interchangeable; incomplete or failed rows cannot be promoted to effects or sensitivity magnitudes, and producing environment lineage remains partial. \\
Permutation checks & \emph{Artifact-supported}: validated CausalForestDML and LinearDML records contain 10 trials with seed 42 for both treatment-column and outcome permutations in each archived dataset--exposure--sampling combination.  Treatment permutations alter the analyzed proxy column; outcome permutations alter mortality in a copied outcome object.  \emph{Not applicable}: CausalPFN permutations were intentionally skipped. & These checks diagnose behavior under disruption; they are not randomization tests or identification proof.  Temporary trial artifacts were removed after aggregation, so only checked aggregate records remain. \\
Producing lineage & \emph{Artifact-supported}: causal run summaries retain stage argv arrays, statuses, return codes, and limited interpreter paths; predictive logs retain argument records and limited device fields; result and reproducibility packages retain hashes and schemas. & \emph{Unresolved}: missing processed-pickle hashes and split manifests; incomplete predictive producing commands and checkpoint mapping; missing numbered causal configs; incomplete source, package, hardware, and archive-copy lineage.  Requirements are intended dependencies, not producing-environment proof. \\
\end{longtable}
\endgroup

\section{Matching and Empirical Support}
\label{sec:robustness-sensitivity-validation-design:overlap-support-diagnostics}

Matching provides a secondary descriptive comparison and indirect support evidence under its implemented binary representation.  For each proxy exposure, the summaries record available pairs, match rate, final Hamming distance, a sufficient-pair flag, maximum-distance warnings, and explicit failures.  Binary confounders are retained, while other numeric confounders may be median-binarized; greedy one-to-one matching then searches within an allowed distance.  A larger accepted distance therefore indicates weaker exact similarity in this representation, not a calibrated distance in the original covariate space.  Failure to obtain pairs is informative support evidence rather than a zero effect.  Matched outcome differences remain descriptive and estimator-dependent; successful matching or a high match rate does not prove positivity, identification, or balance.

Dedicated propensity distributions, common-support plots, pre/post-match balance tables, standardized mean differences, and weighting effective sample sizes were not present in the approved archive.  The available matching fields help identify unsupported comparisons but do not establish global or local positivity, covariate balance, exchangeability, or transportability \cite{crump_et_al_2009_limited_overlap}.

\section{Permutation or Disruption Checks}
\label{sec:robustness-sensitivity-validation-design:permutation-checks-reproducibility-validation}

Permutation checks disrupt one relationship while retaining the remaining run structure.  Treatment permutations shuffle only the currently analyzed proxy column; outcome permutations shuffle only the mortality column in a copied outcome object.  The other proxy columns, time-series data, identifiers, graph, configuration, estimator family, and aggregation fields remain fixed as appropriate.  Subprocess and schema failures are recorded rather than silently aggregated, and temporary trial artifacts are removed after summary extraction.  The validated DML records contain ten trials with seed 42 for each dataset--exposure--sampling combination.  These are negative-control-like sanity checks, not formal randomization tests, p-values, proof of calibration, or evidence that the DAG and exposure are valid \cite{sharma_kiciman_2020_dowhy}.  CausalPFN permutations are intentionally absent.

A response to disrupted exposure or outcome labels tests whether the fitted pipeline distinguishes the observed structure from deliberately broken structure under the implemented procedure.  Even a clear response is a sanity check on pipeline behavior, not causal validation or proof that the original estimate is correctly identified.

\section{Omitted-Variable Sensitivity}
\label{sec:robustness-sensitivity-validation-design:sensitivity-robustness-values}

The DML workflow distinguishes method availability, values called during fitting, saved-training direct values, later saved-artifact recomputation, explicitly labelled fallback reconstruction, benchmark-confounder comparisons, and rendered contours.  These evidence types are not numerically interchangeable.  A value must retain its source and status; unavailable, partial, failed, and not applicable are statuses rather than sensitivity magnitudes.  Benchmarking against an observed covariate supplies an interpretable scale but does not bound every possible unmeasured confounder.  Robustness-value and omitted-variable-bias methods motivate the analysis without validating the local artifacts or identification assumptions \cite{cinelli_hazlett_2020_sensitivity,chernozhukov_et_al_2026_ovb_causal_ml,oprescu_et_al_2019_econml}.  No sensitivity contour is admitted as main-text numerical evidence, and CausalPFN lacks an equivalent archived sensitivity family.

This axis asks how conclusions change under a specified model of unobserved confounding.  It does not establish the actual amount of confounding, prove that no confounding exists below a threshold, validate the DAG, or determine whether an estimate is true or false.

\section{Population Perturbation}

The original-versus-outcome-downsampled design studies dependence on population composition and outcome prevalence.  The original cohorts remain primary, and the sampled populations remain separate: downsampling changes the empirical population, nuisance-fitting problem, support, and averaging distribution.  Stability or change across those populations is evidence about this specified perturbation, not a replacement for external validation and not evidence that the original cohort is unbiased.

\section{Provenance, Evidence Admission, and Reproducibility Boundaries}
\label{sec:experimental-design:evaluation-result-selection-policy}

Before a value enters Chapter~\ref{chap:results}, its producing run, dataset, model or estimator, sampling condition, source path, schema, status, and population must be identified.  Failed artifacts are excluded; fallback diagnostics retain their labels; source fields are interpreted at their actual granularity; and the estimand and unresolved limitations remain explicit.  Directory names, copied summaries, or non-null cells alone are insufficient.  This policy preserves the original-cohort and estimator hierarchy and prevents numerical availability from being mistaken for scientific admissibility.

\label{sec:experimental-design:computational-environment-reproducibility}
The result manifest and compact reproducibility package record available paths, hashes, schemas, commands and log arguments, seeds, environment fields, and artifact roles.  They distinguish current source behavior, archived numerical artifacts, checked result rows, reconstructed summaries, unresolved producing lineage, and full clean-checkout reproducibility.  These classes support numerical transcription and execution traceability, but not clean-checkout reproduction or scientific identification.  Raw and processed data and numbered final causal configurations are absent; predictive split and checkpoint lineage, source versions, environments, and archive history are incomplete; and historical paths are not portable.  Requirements and current defaults describe interfaces, not necessarily the producing environment.  Provenance controls what may be claimed from a testbed; provenance quality does not itself establish clinical or causal validity.

\section{Synthesis of the Non-Equivalent Axes}

Each axis asks a different question and supplies a bounded form of evidence.  Construct plausibility asks whether a representation and rationale are inspectable, but cannot establish clinical construct validity.  Proxy recoverability asks whether rule-derived labels are learnable, but cannot establish diagnostic or causal correctness.  Mortality prediction asks whether the representation contains prognostic information, but cannot establish causal relevance.  Cross-estimator sign, rank, and magnitude stability asks whether a pattern depends on one fitted estimator, but cannot establish accuracy, common estimands, or truth.  Clinical-literature comparison asks whether patterns are broadly compatible with heterogeneous external studies, but cannot provide a meta-analysis, external proxy validation, or calibration against causal truth.

Matching and empirical-support diagnostics ask whether usable comparisons are available under the implemented representation, but cannot prove positivity, balance, exchangeability, or identification.  Permutation or disruption checks ask whether deliberately broken labels or outcomes change fitted behavior, but cannot validate the original causal analysis.  Omitted-variable sensitivity asks how conclusions respond to a specified confounding model, but cannot reveal the actual confounding strength or validate the DAG.  Population perturbation asks whether findings depend on the original-versus-downsampled population, but cannot establish external validity or remove selection bias.  Provenance and evidence admission ask which artifacts and claims are traceable, but cannot establish clinical meaning or causal correctness.  Together these axes characterize complementary strengths, discrepancies, and failure modes without producing an evaluation score or a causal answer key.
